\section{Related Work}
\label{sec:related_work}

Our work is situated at the intersection of weight-space model merging, dynamic routing, and parameter-efficient multi-task adaptation. 

\paragraph{Static Weight Merging}
Model merging aims to combine multiple independently fine-tuned models sharing a common pre-trained initialization without further training. Early approaches relied on direct linear interpolation of weights, such as Model Soups \cite{wortsman2022model} and Federated Averaging \cite{mcmahan2017communication}. While highly effective for models trained on similar tasks, direct averaging can suffer from representation misalignment due to permutation symmetries. To mitigate this, coordinate alignment methods like Git Re-Basin \cite{ainsworth2022git} align neural activations before averaging. More recently, task arithmetic \cite{task_arithmetic} demonstrates that adding "task vectors" (the difference between fine-tuned and pre-trained weights) can steer models toward specific behaviors. To prevent destructive interference from overlapping weight edits in task arithmetic, TIES-Merging \cite{yadav2023ties} prunes small parameter changes and resolves sign disagreements. Similarly, DARE \cite{yu2023dare} uses random dropping and rescaling of weight edits, and Fisher Merging \cite{matena2022merging} weights parameter combinations using the diagonal Fisher information matrix. However, all these methods remain static and task-agnostic at runtime, performing poorly on highly heterogeneous stream deployment.

\paragraph{PEFT Weight Merging and LoRA}
Our work is also closely related to Parameter-Efficient Fine-Tuning (PEFT) techniques, particularly Low-Rank Adaptation (LoRA) \cite{hu2021lora}, which fine-tunes small low-rank update matrices instead of entire model checkpoints. Merging multiple LoRA modules dynamically at runtime (e.g., LoRA-MoE or LoraHub) has emerged as a lightweight paradigm for multi-task adaptation. Since BWS-Router operates on task vectors (the differences between fine-tuned and pre-trained weights), it can be seamlessly applied to dynamically blend low-rank LoRA matrices. Under this regime, BWS-Router's block-wise parameter sharing would yield even greater relative parameter and routing compute savings, making it highly compatible with modern PEFT deployment pipelines.

\paragraph{Dynamic Routing and Mixture of Experts}
To handle diverse runtime streams, Mixture of Experts (MoE) \cite{jacobs1991adaptive, shazeer2017outrageously, switch_transformer} employs input-conditioned routing gates to dynamically select a sparse subset of expert subnetworks. While highly successful in large language models, MoE is computationally demanding because it routes token-level activations through separate feed-forward layers, requiring specialized distributed infrastructure and significant memory footprints during inference. Dynamic model merging, on the other hand, performs routing directly in the parameter space (weight space) rather than activation space. It dynamically blends the entire parameter matrices of expert networks on a per-sample or per-batch basis, introducing zero inference latency and requiring no architectural changes to the backbone.

\paragraph{Progression of Dynamic Model Merging}
The evolution of dynamic model merging in the literature reveals a persistent tension between gating complexity and generalization robustness. Early formulations, such as Routing Soups \cite{routing_soup} and the Bounded Classical Router (\textbf{BC-Router}) family (comprising BSigmoid-Router and GLS-Router) \cite{bc_router}, analyzed the behavior of the standard Softmax gating activation. While Softmax gating is highly optimal and empirically superior for closed-world classification tasks (where classes are mutually exclusive and implicit sum-to-one regularizing constraints prevent parameter scaling drift), it acts as a zero-sum competitive bottleneck that cannot handle decoupled open-world streams, where we must co-activate multiple experts simultaneously or deactivate them entirely under out-of-distribution (OOD) inputs. To address this, the Layer-wise Low-dimensional Classical Router (\textbf{L3-Router}) \cite{l3_router_orig} introduced unshared, layer-wise independent routers mapping unsupervised low-dimensional PCA projections onto task coefficients. 

To introduce non-monotonic gating dynamics, Quantum Wavefunction Superposition Merging (\textbf{QWS-Merge}) \cite{qws_merge_orig} formulated wave-inspired routing equations utilizing complex trigonometric amplitudes. While theoretically elegant, we show that QWS-Merge's highly non-linear activation landscape is extremely non-convex and rugged, leading to severe optimization instability and frequent training collapse across random seeds. Furthermore, we identify that unshared architectures like L3-Router suffer from high-frequency coefficient ruggedness, parameter scaling excess, and heavy computational overhead on small calibration datasets (64 samples). Our proposed \textbf{BWS-Router} addresses these exact issues. By introducing block-wise weight sharing, we constrain the optimization search space, reduce parameter counts by 66.7\%, and mathematically prove the mitigation of layer-to-layer coefficient ruggedness, establishing a stable, parameter-efficient, and robust dynamic model merging framework.
